\documentclass[11pt,a4paper]{article}

% Minimal authoring template for daily writing.
% Keep this lean; integrate venue-specific formatting only in target-repository submission assets when needed.
% Rule: references shall be bound to the preceding text using '~' (e.g., Figure~\ref{fig:example}, see~\cite{Key2024}).

\usepackage[utf8]{inputenc}
\usepackage[T1]{fontenc}
\usepackage[english]{babel}
\usepackage{lmodern}
\usepackage{microtype}
\usepackage[a4paper,margin=2.5cm]{geometry}
\usepackage{setspace}
\usepackage{graphicx}
\usepackage{booktabs}
\usepackage{amsmath,amssymb}
\usepackage{authblk}
\usepackage[numbers,sort&compress]{natbib}
\usepackage[hidelinks]{hyperref}

\onehalfspacing

\title{Paper Title (EN)}
\author[1]{Hugo Ormo}
\author[2]{Hugo Ormo}
\affil[1]{GfSE UAF Working Group}
\affil[2]{NTT DATA Deutschland SE, 80807 M\"unchen, \href{mailto:hugo.ormo@nttdata.com}{hugo.ormo@nttdata.com}}
\date{}

\begin{document}
\maketitle

\textbf{Practical contribution:} TBD


\begin{abstract}
Replace with your abstract. Cite early while drafting, e.g. \cite{Weilkiens2024_SysMLv2Book}.
\end{abstract}

\textbf{Keywords:} SysML v2, MBSE, architecture

\section{Introduction}
Start writing directly in LaTeX.

\section{Approach}
Add your method and evidence.

\section{Discussion}
Interpret results and limitations.

\section{Conclusion}
Summarize key contributions and next steps.

\section*{Acknowledgments}
Artificial intelligence tools were used in the elaboration of this document, including support for drafting, refinement, and editorial improvement.

\bibliographystyle{IEEEtran}
\bibliography{references}

\end{document}
